\chapter{Modelos Lineares Generalizados}
\label{cap:glm}
% ── índice ──────────────────────────────────────────────────────
\index{modelos lineares generalizados|textbf}
\index{GLM|see{modelos lineares generalizados}}
\index{glm@\texttt{glm()}|textbf}
\index{PPML|textbf}
\index{função de ligação|textbf}
\index{gamma, distribuição}
\index{equações de estimação generalizadas}

% ─────────────────────────────────────────────────────────────────────────────
\section{A família exponencial}

Os Modelos Lineares Generalizados (Nelder e Wedderburn 1972) unificam OLS,
Logit, Poisson e muitos outros estimadores sob um arcabouço comum. Três
componentes definem um GLM:

\begin{enumerate}
  \item \textbf{Componente aleatório}: $Y_i \sim$ família exponencial com
  média $\mu_i$.
  \item \textbf{Preditor linear}: $\eta_i = x_i'\beta$.
  \item \textbf{Função de ligação}: $g(\mu_i) = \eta_i$, que conecta média
  à escala linear.
\end{enumerate}

\begin{center}
\begin{tabular}{llll}
\toprule
\textbf{Família} & \textbf{Link padrão} & \textbf{Variância} & \textbf{Aplicação} \\
\midrule
Gaussiana   & identidade & $\sigma^2$   & OLS \\
Binomial    & logit      & $\mu(1-\mu)$ & dados binários \\
Poisson     & log        & $\mu$        & contagens \\
Gamma       & log        & $\mu^2/\phi$ & dados positivos com caudas \\
\bottomrule
\end{tabular}
\end{center}

A estimação por IRLS (Iteratively Reweighted Least Squares) é eficiente e
converge rapidamente.

\subsection*{GLM Gamma — dados positivos com caudas pesadas}

Gastos, tamanho de empréstimos e durações seguem distribuições assimétricas
à direita com variância proporcional ao quadrado da média. O GLM Gamma com
link log modela $E[Y \mid x] = \exp(x'\beta)$ com $\mathrm{Var}[Y\mid x] \propto \mu^2$.

\subsection*{PPML — comércio internacional}

O estimador PPML (Poisson Pseudo MLE) de Santos Silva e Tenreyro (2006) usa
a equação de estimação Poisson mesmo quando $Y$ não é inteira (ex: valores
de comércio). Ele é robusto a zeros e a heterocedasticidade multiplicativa
— problemas comuns em equações gravitacionais de comércio.

% ─────────────────────────────────────────────────────────────────────────────
\section{Verificação por DGP}

\begin{lstlisting}[caption={GLM Gamma e PPML}]
seed(42)
generate df x  = rnormal()
generate df mu = exp(0.5 + 1.0 * x)
// Gamma: média=mu, variância ~ mu² (simulado via lognormal)
generate df y_gamma = mu * exp((rnormal() - 0.5) * 0.6)

let m_gamma = glm(y_gamma ~ x, df, family=gamma, link=log)
display m_gamma

// PPML com dados de "comércio" (mistura zeros e positivos)
generate df x2    = rnormal()
generate df trade = round(exp(0.3 + 1.2*x + 0.8*x2) * (rnormal() > 0.2))
let m_ppml = glm(trade ~ x + x2, df, family=poisson, link=log)
display m_ppml
\end{lstlisting}

\subsection*{GLM Gamma}

\begin{lstlisting}[language={}, basicstyle=\ttfamily\small,
  frame=none, numbers=none, backgroundcolor=\color{codbg},
  xleftmargin=0pt, xrightmargin=0pt, breaklines=false]
====================== Generalized Linear Model Results ======================
Family: Gamma  Link: Log  N=500  Dispersion=0.0297
------------------------------------------------------------------------------
Variable    coef    std err      z    P>|z|
const      0.483     0.015   31.49    0.000    ***
x          1.065     0.027   39.03    0.000    ***
\end{lstlisting}

$\hat\beta_0 = 0{,}483 \approx 0{,}5$ e $\hat\beta_1 = 1{,}065 \approx 1{,}0$.
A dispersão $\hat\phi = 0{,}030$ indica variância relativa pequena.

\subsection*{PPML}

\begin{lstlisting}[language={}, basicstyle=\ttfamily\small,
  frame=none, numbers=none, backgroundcolor=\color{codbg},
  xleftmargin=0pt, xrightmargin=0pt, breaklines=false]
Family: Poisson  Link: Log  N=500  Pseudo R²=0.2936
------------------------------------------------------------------------------
Variable    coef    std err      z    P>|z|
const      0.069     0.077    0.91    0.366
x          1.168     0.090   12.93    0.000    ***
x2         0.890     0.090    9.91    0.000    ***
\end{lstlisting}

PPML recupera $\hat\beta_x \approx 1{,}17$ (verdadeiro: 1{,}2) e
$\hat\beta_{x_2} \approx 0{,}89$ (verdadeiro: 0{,}8), lidando com os zeros
da variável \texttt{trade}.

% ─────────────────────────────────────────────────────────────────────────────
\section{Sintaxe}

\begin{lstlisting}
glm(y ~ x, df, family=gaussian, link=identity)  // OLS via GLM
glm(y ~ x, df, family=binomial, link=logit)     // Logit via GLM
glm(y ~ x, df, family=gamma,    link=log)       // Gamma com link log
glm(y ~ x, df, family=poisson,  link=log)       // Poisson / PPML
\end{lstlisting}

\begin{center}
\begin{tabular}{ll}
\toprule
\textbf{Família} & \textbf{Links disponíveis} \\
\midrule
\texttt{gaussian}  & \texttt{identity}, \texttt{log}, \texttt{inverse} \\
\texttt{binomial}  & \texttt{logit}, \texttt{probit}, \texttt{cloglog} \\
\texttt{poisson}   & \texttt{log}, \texttt{identity}, \texttt{sqrt} \\
\texttt{gamma}     & \texttt{log}, \texttt{identity}, \texttt{inverse} \\
\bottomrule
\end{tabular}
\end{center}
